\documentclass[12pt,a4paper]{article}

% ── Packages ────────────────────────────────────────────────────────────────
\usepackage[top=0.9in, bottom=0.9in, left=1in, right=1in]{geometry}
\usepackage{setspace}
\usepackage{booktabs}
\usepackage{array}
\usepackage{amsmath}
\usepackage{graphicx}
\usepackage{float}
\usepackage{caption}
\usepackage[hidelinks]{hyperref}
\usepackage{microtype}
\usepackage{titlesec}
\usepackage{natbib}
\usepackage{threeparttable}
\usepackage{etoolbox}
\apptocmd{\thebibliography}{\small\setlength{\itemsep}{1pt}\setlength{\parskip}{0pt}}{}{}

% ── Typography ───────────────────────────────────────────────────────────────
\setstretch{1.08}
\setlength{\parindent}{1.2em}
\setlength{\parskip}{0pt}

\titleformat{\section}{\normalsize\bfseries\scshape}{}{0em}{}
\titleformat{\subsection}{\normalsize\bfseries}{}{0em}{}
\titleformat{\paragraph}[runin]{\normalsize\bfseries}{}{0em}{}[.]
\titlespacing*{\section}{0pt}{7pt}{2pt}
\titlespacing*{\subsection}{0pt}{5pt}{2pt}
\titlespacing*{\paragraph}{0pt}{3pt}{5pt}

\captionsetup{font=small, skip=2pt}

% ── Title ────────────────────────────────────────────────────────────────────
\title{\vspace{-2em}\large\bfseries Does AI Compress Wages?\\[0.2em]
\normalsize Evidence from a Trend-Break Difference-in-Differences Design}
\author{\small Yishu Sheng \quad Shanqi Zhang \quad Leyla Aysegul Kandemir\\[1pt]
\small University of Chicago Booth School of Business \quad
BUS 41207: Causal Inference \quad March 2026}
\date{}

\begin{document}
\maketitle
\vspace{-2em}

% ── Abstract ─────────────────────────────────────────────────────────────────
\begin{center}
\begin{minipage}{0.9\textwidth}
\footnotesize\setlength{\parindent}{0pt}
\textbf{Abstract.}
Empirical studies of generative AI's labor-market effects reach conflicting conclusions despite using similar designs. We show this divergence reflects an identification problem: occupations with higher AI exposure followed different wage trajectories before ChatGPT, violating the parallel-trends assumption in all 12 placebo tests across four major exposure indices. Using the trend-break estimator of \citet{wolfers2006}, we find that ChatGPT \emph{raised} wages by $+0.31\%$ in high-exposure occupations — reversing the sign of standard DID — while leaving aggregate employment unchanged. Wage gains are concentrated at the bottom of the distribution ($+0.62\%^{***}$ at $p_{10}$, zero at $p_{90}$), consistent with task upgrading. A firm-level extension finds selective headcount reductions at more exposed public firms, suggesting restructuring coexists with stable aggregate employment.
\end{minipage}
\end{center}

\vspace{4pt}

% ── 1. Introduction / Background ─────────────────────────────────────────────
\section{Introduction}

Does generative AI raise wages, lower them, or leave them unchanged? The question matters for compensation strategy, workforce planning, and labor policy. Yet empirical estimates diverge sharply: studies using similar difference-in-differences (DID) designs around ChatGPT's November 2022 launch report positive, negative, and null wage effects for similar occupational groups \citep{acemoglu2022, brynjolfsson2025, hui2024}. Causal inference is essential because occupations and firms with high AI exposure differ systematically from those with low exposure even \emph{before} ChatGPT — they tend to be higher-skill, higher-wage, and already on different growth paths \citep{alm2003}. A naive pre–post comparison risks attributing these pre-existing differences to AI. This paper diagnoses and corrects that identification problem.

Our contribution is threefold. First, we demonstrate through 12 placebo tests that the standard parallel-trends assumption fails for every major AI exposure index currently used in the literature. Second, we adopt the trend-break estimator of \citet{wolfers2006} to absorb pre-existing differential trends and identify the post-ChatGPT causal break. Third, we combine occupation-level BLS evidence with a firm-level Compustat extension that passes its own pre-trend checks. After correction, the wage effect reverses sign and the employment effect disappears, pointing to wage compression via task upgrading rather than displacement.

% ── 2. Literature Review ─────────────────────────────────────────────────────
\section{Related Literature}

The task-based framework of \citet{alm2003} predicts that computers substitute for routine tasks while complementing non-routine cognitive work. \citet{acemoglu2018} extend this, showing the labor share falls when automation outpaces new task creation. Our floor-raiser result contrasts with both frameworks: generative AI appears to raise wages most for routine and manual workers.

On exposure measurement, \citet{felten2021} map AI benchmarks onto O*NET tasks; \citet{eloundou2024} extend this to GPT capabilities; \citet{webb2020} use patent similarity; \citet{fo2013} provide an earlier automation-risk baseline. We show all four produce biased DID estimates in our setting. On empirical AI effects, \citet{brynjolfsson2025} and \citet{noy2023} find productivity gains concentrated among less-experienced workers. \citet{hui2024} find employment losses on freelance platforms.

Methodologically, we build on \citet{wolfers2006} and the recent DID literature. \citet{roth2023} survey violations of parallel trends; \citet{callaway2021} and \citet{sun2021} address staggered adoption, which is not our setting — our challenge is exposure-correlated pre-trends with a common treatment date.

% ── 3. Data ──────────────────────────────────────────────────────────────────
\section{Data}

\paragraph{BLS OEWS} We use the BLS Occupational Employment and Wage Statistics for 2019--2024, compiled into a state $\times$ occupation $\times$ year panel of 157,677 observations covering 667 occupations and 51 states. Outcomes are log median annual wage, log employment, and log wages at the 10th, 25th, 75th, and 90th percentiles. We drop observations with BLS-suppressed wage values (approximately 3,500 cells at the 90th percentile).

\paragraph{AI exposure indices} Our baseline treatment is the AIOE index \citep{felten2021}, standardized to mean zero and unit standard deviation. We also construct GPT exposure \citep{eloundou2024} from original task-level data using O*NET core task weights, Frey-Osborne automation risk \citep{fo2013}, and Webb's patent-based exposure index \citep{webb2020}. All four are merged to the panel as time-invariant occupation-level measures. The correlation between AIOE and GPT exposure is $r=0.90$; Frey-Osborne correlates negatively with AIOE ($r=-0.42$), reflecting its focus on industrial-era substitution rather than LLM capability.

\paragraph{O*NET task characteristics} We merge O*NET routine task intensity, manual ability, computer use, and cognitive ability scores for the mechanism analysis.

\paragraph{Compustat} For the firm-level extension, we use Compustat North America for 2019--2024, restricting to USD-denominated industrial firms with non-missing employment and total assets, yielding 5,417 firms and 24,845 firm-year observations. The key treatment variable is an employment-weighted Bartik shift-share using the Anthropic Economic Index \citep{massenkoff2026}:
\[
\text{AEI Exposure}_j = \sum_k \frac{\text{emp}_{jk}}{\sum_k \text{emp}_{jk}} \times \text{observed\_exposure}_k
\]
aggregated across 231 4-digit NAICS industries. Top industries: Business Support (0.43), Securities (0.36), Software Publishers (0.32).

% ── 4. EDA ───────────────────────────────────────────────────────────────────
\section{Exploratory Data Analysis}

\paragraph{Distribution of AI exposure} The AIOE index has a roughly symmetric distribution across occupations (mean 3.5, range 1.4--6.5 in our sample). High-exposure occupations include lawyers (0.84), data analysts (0.79), and accountants (0.73); low-exposure occupations include plumbers (0.18), truck drivers (0.22), and janitors (0.11). The GPT Exposure index is highly correlated with AIOE ($r=0.90$), while Frey-Osborne is negatively correlated ($r=-0.42$) because it captures industrial-era substitution risk rather than LLM complementarity. Webb is near-orthogonal ($r=0.05$), measuring patent-based technological exposure. These differences across indices motivate our multi-index placebo strategy: if the bias is structural and not index-specific, all four should fail.

\paragraph{Pre-treatment wage trends} Figure~\ref{fig:pretreats} plots mean log wages for occupations above and below the median AIOE score from 2019 to 2024. Two features stand out: high-AIOE occupations have persistently higher wages, and the gap between the two groups narrows steadily during 2019--2022, \emph{before} ChatGPT. This pre-treatment convergence directly violates the parallel-trends assumption and motivates the identification correction in Section~\ref{sec:methods}. The firm-level data tells a different story: plotting employment trends for high- versus low-AEI-exposure industries shows no systematic divergence during 2019--2021, confirming that the Bartik instrument satisfies parallel trends at the firm level even where the occupation-level indices do not.

\begin{figure}[H]
\centering
\includegraphics[width=0.75\textwidth]{figure_1.pdf}
\caption{Log median wage by AIOE exposure group, 2019--2024. The pre-treatment convergence (2019--2022) predates ChatGPT and violates parallel trends.}
\label{fig:pretreats}
\end{figure}

% ── 5. Methods ───────────────────────────────────────────────────────────────
\section{Empirical Strategy}
\label{sec:methods}

\subsection{Standard DID and Placebo Diagnosis}

The standard specification is:
\begin{equation}
Y_{sot} = \alpha_{os} + \gamma_t + \phi_{st} + \beta(\text{AIOE}_o \times \text{Post}_t) + \varepsilon_{sot}
\label{eq:std}
\end{equation}
with occupation$\times$state, year, and state$\times$year fixed effects. Standard errors are clustered by occupation$\times$state. Applying \eqref{eq:std} gives $\hat\beta=-2.01\%^{***}$ for wages — a large negative effect. We test whether this is spurious by substituting fake treatment dates of 2020, 2021, and 2022 across all four exposure indices.

\begin{table}[H]
\centering
\caption{Placebo tests: all 12 cases reject $p<0.001$ before treatment}
\label{tab:placebo}
\small
\begin{tabular}{lcccc}
\toprule
\textbf{Index} & \textbf{Fake 2020} & \textbf{Fake 2021} & \textbf{Fake 2022} & \textbf{True 2023} \\
\midrule
AIOE     & $-0.0208^{***}$ & $-0.0224^{***}$ & $-0.0222^{***}$ & $-0.0201^{***}$ \\
GPT Exp. & $-0.0201^{***}$ & $-0.0223^{***}$ & $-0.0217^{***}$ & $-0.0196^{***}$ \\
F\&O     & $+0.0124^{***}$ & $+0.0138^{***}$ & $+0.0138^{***}$ & $+0.0126^{***}$ \\
Webb     & $-0.0054^{***}$ & $-0.0070^{***}$ & $-0.0058^{***}$ & $-0.0044^{***}$ \\
\midrule
$N$ obs  & \multicolumn{4}{c}{157,677} \\
\bottomrule
\end{tabular}
\begin{flushleft}
\footnotesize FE: Occ$\times$State + Year + State$\times$Year. SE clustered by occ$\times$state. The bias is structural — it holds regardless of which exposure index is used.
\end{flushleft}
\end{table}

The smoking gun is clearest when we substitute routine task intensity — a pre-AI occupational characteristic — as the treatment: standard DID gives $+0.83\%^{***}$, collapsing to $-0.07\%$ (n.s.) under the trend-break correction. This confirms that standard DID conflates pre-existing polarization dynamics with the causal effect of ChatGPT.

\subsection{Trend-Break DID}

We adopt the \citet{wolfers2006} trend-break estimator:
\begin{equation}
Y_{sot} = \alpha_{os} + \gamma_t + \phi_{st} + \delta_1(\text{AIOE}_o \times t) + \delta_2(\text{AIOE}_o \times \text{Post}_t) + \varepsilon_{sot}
\label{eq:tb}
\end{equation}
$\delta_1$ absorbs the pre-existing differential trend; $\delta_2$ is the causal break at ChatGPT's launch. The identifying assumption — that the pre-existing trend would have continued absent treatment — is strictly weaker than parallel trends.

\paragraph{Threats to validity} The design cannot address non-linear confounders; we check with quadratic and cubic trend specifications. Estimates are relative effects by exposure intensity, not treatment-versus-control contrasts. The post-period is short, so results capture near-term rather than long-run effects. Finally, the firm-level analysis covers only publicly listed companies and may not generalize to small and medium enterprises.

% ── 6. Results ───────────────────────────────────────────────────────────────
\section{Results}

\subsection{Main Estimates and Sign Reversal}

\begin{table}[H]
\centering
\caption{Main results: sign reversal from trend-break correction}
\label{tab:main}
\small
\begin{tabular}{lcc}
\toprule
 & \textbf{Log Median Wage} & \textbf{Log Employment} \\
\midrule
Standard DID: $\hat\beta$ (AIOE $\times$ Post)     & $-0.0201^{***}$ & $+0.0211^{***}$ \\[3pt]
Trend-Break: $\hat\delta_1$ (AIOE $\times t$)       & $-0.0078^{***}$ & $+0.0077^{***}$ \\
Trend-Break: $\hat\delta_2$ (AIOE $\times$ Post)    & $+0.0031^{***}$ & $-0.0018$ \\
\quad $p$-value                                      & $<0.001$        & $0.495$ \\
\midrule
$N$ obs & \multicolumn{2}{c}{157,677} \\
\bottomrule
\end{tabular}
\begin{flushleft}
\footnotesize All specifications include Occ$\times$State + Year + State$\times$Year FE. SE clustered by occ$\times$state. $^{***}p<0.001$.
\end{flushleft}
\end{table}

The central result is the sign reversal: standard DID implies $-2.01\%^{***}$; the trend-break correction gives $+0.31\%^{***}$. The pre-trend $\hat\delta_1=-0.0078^{***}$ confirms that high-AIOE occupations were already converging toward lower-exposure occupations before ChatGPT — this is what standard DID was picking up. Once absorbed, ChatGPT \emph{raised} wages in high-exposure occupations. The employment effect is indistinguishable from zero at the occupation level.

\subsection{Distributional Effects: Wage Floor}

\begin{table}[H]
\centering
\caption{Trend-break estimates by wage percentile ($\hat\delta_2$, AIOE $\times$ Post)}
\label{tab:pctile}
\small
\begin{tabular}{lccccc}
\toprule
 & $p_{10}$ & $p_{25}$ & $p_{50}$ & $p_{75}$ & $p_{90}$ \\
\midrule
$\hat\delta_2$ & $+0.0062^{***}$ & $+0.0055^{***}$ & $+0.0031^{***}$ & $+0.0024^{**}$ & $+0.0002$ \\
SE             & $(0.0011)$      & $(0.0009)$      & $(0.0008)$      & $(0.0009)$     & $(0.0010)$ \\
\bottomrule
\end{tabular}
\begin{flushleft}
\footnotesize SE clustered by occ$\times$state. $^{***}p<0.001$, $^{**}p<0.01$. Effect is monotonically declining from $p_{10}$ to $p_{90}$.
\end{flushleft}
\end{table}

Figure~\ref{fig:percentiles} and Table~\ref{tab:pctile} confirm a monotonic pattern: the gain is $+0.62\%^{***}$ at the 10th percentile and statistically indistinguishable from zero at the 90th. The $p_{90}/p_{10}$ ratio falls by $-0.54\%^{***}$, confirming wage compression. This is inconsistent with a skill-complementarity story, which predicts gains concentrated at the top.

\begin{figure}[H]
\centering
\includegraphics[width=0.65\textwidth]{figure_2.pdf}
\caption{Trend-break wage effects ($\hat\delta_2$) by wage percentile with 95\% CI. Effect declines monotonically — AI raises the wage floor.}
\label{fig:percentiles}
\end{figure}

\subsection{Mechanism: Task Upgrading}

To distinguish task upgrading from skill complementarity, we split the sample at the median of each O*NET task dimension and estimate the trend-break DID separately. Skill complementarity predicts gains concentrated in cognitive and computer-intensive occupations; task upgrading predicts gains in routine and manual occupations where AI automates the repetitive components of jobs.

\begin{table}[H]
\centering
\caption{Heterogeneity by occupation type: $\hat\delta_2$ (AIOE $\times$ Post), log median wage}
\label{tab:het}
\small
\begin{tabular}{lcccc}
\toprule
\textbf{O*NET Dimension} & \textbf{High $\hat\delta_2$} & \textbf{High $N$} & \textbf{Low $\hat\delta_2$} & \textbf{Low $N$} \\
\midrule
Routine Tasks     & $+0.0057^{**}$  & 76,140 & $+0.0046^{***}$ & 81,537 \\
SE                & $(0.0019)$      &        & $(0.0014)$      & \\[3pt]
Manual Ability    & $+0.0055^{*}$   & 74,100 & $+0.0013$       & 83,577 \\
SE                & $(0.0026)$      &        & $(0.0018)$      & \\[3pt]
Computer Use      & $-0.0014$       & 81,952 & $+0.0028$       & 75,725 \\
SE                & $(0.0017)$      &        & $(0.0018)$      & \\[3pt]
Cognitive Ability & $-0.0016$       & 81,219 & $+0.0037^{*}$   & 76,458 \\
SE                & $(0.0017)$      &        & $(0.0016)$      & \\
\bottomrule
\end{tabular}
\begin{flushleft}
\footnotesize FE: Occ$\times$State + Year + State$\times$Year. SE clustered by occ$\times$state. $^{***}p<0.001$, $^{**}p<0.01$, $^{*}p<0.05$. Split at above/below median O*NET score.
\end{flushleft}
\end{table}

The pattern is inconsistent with skill complementarity. Cognitive ($-0.16\%$, n.s.) and computer-intensive ($-0.14\%$, n.s.) workers — the groups that standard skill-complementarity theory predicts should benefit most — show no significant wage gains. Instead, the largest and most precisely estimated gains accrue to routine task workers ($+0.57\%^{**}$) and manual ability workers ($+0.55\%^{*}$). Representative winners: data entry clerks, billing specialists, payroll workers, bank tellers, food prep workers, assemblers. Representative non-winners: lawyers, physicians, software engineers, financial analysts. This is consistent with AI automating the repetitive, rule-based components of lower-paid jobs, freeing workers to perform higher-value tasks without requiring skill changes — a task-upgrading story.

\subsection{Firm-Level Extension}

\begin{table}[H]
\centering
\caption{Firm-level DID: Compustat 2019--2024 (AEI Bartik treatment)}
\label{tab:firm}
\small
\begin{tabular}{lcccc}
\toprule
\textbf{Outcome} & \textbf{Coef.} & \textbf{SE} & \textbf{Sig.} & \textbf{$N$} \\
\midrule
$\log(\text{Employment})$            & $-0.263$ & $(0.088)$ & $^{**}$ & 24,845 \\
$\log(\text{Labor Cost}/\text{emp})$ & $+0.030$ & $(0.152)$ & n.s. & 4,729 \\
$\log(\text{Revenue})$               & $-0.168$ & $(0.112)$ & n.s. & 22,299 \\
\bottomrule
\end{tabular}
\begin{flushleft}
\footnotesize Firm + Year FE. SE clustered by firm. Control: $\log(\text{assets})$. $^{**}p<0.01$. Pre-trend test passes (placebo $-0.114$, n.s.). Labor cost missingness (81\%) reflects voluntary disclosure of \texttt{xlr}.
\end{flushleft}
\end{table}

Unlike the occupation-level analysis, the firm-level design passes its pre-trend check: event study coefficients for 2019, 2020, and 2021 are all statistically indistinguishable from zero, and a placebo test using a fake 2022 treatment date produces $-0.114$ (n.s.). Employment falls significantly ($-0.263^{**}$, $p<0.01$): a firm at the 75th vs.\ 25th percentile of industry exposure (exposure scores 0.225 vs.\ 0.105) shows roughly 3 percentage points lower employment growth by 2024. The effect grows over time — $-0.145^{*}$ in 2023, $-0.281^{***}$ in 2024 — consistent with gradual AI adoption rather than a one-time shock.

Labor costs per worker and revenue are unchanged. If labor cost per employee is flat while headcount falls, firms are selectively reducing employment rather than cutting wages across the board. If revenue is also flat, the same output is being produced with fewer workers — a productivity story, not a demand shock. Together this points to firms restructuring their workforce composition in response to AI, not experiencing a contraction in their business. This reconciles the occupation-level null employment result with the firm-level negative: headcount falls at large public firms through slower hiring, not mass layoffs, so workers are reallocated rather than unemployed.

% ── 7. Robustness ────────────────────────────────────────────────────────────
\section{Robustness Checks}

\begin{table}[H]
\centering
\caption{Summary of robustness checks (outcome: log median wage unless noted)}
\label{tab:robust}
\small
\begin{tabular}{p{5.8cm}p{6.5cm}}
\toprule
\textbf{Check} & \textbf{Result} \\
\midrule
Quadratic trend-break        & $\hat\delta_2=+0.0026^{**}$ — holds \\
Cubic trend-break            & $\hat\delta_2=+0.0048^{***}$ — holds, stronger \\
Drop COVID years (2020--21)  & $\hat\delta_2=+0.0040^{***}$ — pandemic attenuated result \\
Exclude tech sector (SOC 15) & $\hat\delta_2=+0.0031^{***}$ — unchanged \\
Balanced panel               & $\hat\delta_2=+0.0026^{**}$ — holds \\
Two-way clustering           & $\hat\delta_2=+0.0031^{*}$, $p=0.014$ — holds \\
Direct inequality test $\log(p_{90}/p_{10})$ & $\hat\delta_2=-0.0054^{***}$ — compression confirmed \\
CPS 72-month event study     & No employment break post-ChatGPT \\
ACS individual-level (190K)  & $\hat\delta_2(p_{10})$ is $5.4\times$ larger than $\hat\delta_2(p_{90})$ \\
\bottomrule
\end{tabular}
\end{table}

The floor-raiser finding survives all checks (Table~\ref{tab:robust}). Dropping COVID years strengthens the wage estimate to $+0.40\%^{***}$, ruling out pandemic confounding. Two-way clustering by occupation and state drops significance from $p<0.001$ to $p=0.014$ but the sign and magnitude are unchanged. A direct inequality test using $\log(p_{90}/p_{10})$ as outcome gives $\hat\delta_2=-0.0054^{***}$, confirming wage compression in a single regression. Notably, the pre-trend for this inequality regression is $\hat\delta_1=+0.0050^{***}$ — inequality was \emph{widening} before ChatGPT, which standard DID would have misattributed as a narrowing treatment effect. The CPS 72-month event study (25,114 obs, 351 occupations) confirms no structural employment break post-ChatGPT. The ACS validation (164,056 obs) is directionally consistent with the floor-raiser: the $p_{10}$ effect is 5.4$\times$ larger than the $p_{90}$ effect, though neither reaches statistical significance at the lower occupation-level (vs.\ occupation$\times$state-level) power of this dataset.

% ── 8. Conclusion ────────────────────────────────────────────────────────────
\section{Conclusion and Business Implications}

Standard DID estimates of AI's labor-market effects are systematically misleading in our setting because AI exposure is correlated with pre-existing wage trends. Once we absorb those trends, ChatGPT raised wages in high-exposure occupations by $+0.31\%^{***}$, with gains concentrated at the bottom of the distribution — a floor-raiser, not a ceiling-raiser. Aggregate employment is unchanged at the occupation level, while firm-level evidence points to selective headcount reductions consistent with productivity restructuring. The mechanism is task upgrading rather than skill complementarity: routine and manual workers gain most, while lawyers and software engineers see no significant effect.

\paragraph{Business implications} Firms that adopt AI to increase productivity will face upward wage pressure at the bottom of the distribution as outside options for routine workers improve. The null labor cost result at the firm level suggests firms are currently capturing the productivity surplus from AI rather than passing it through to wages, but competitive dynamics may shift this as adoption becomes more uniform across industries. Managers should expect AI's effects to show up first through task redesign, wage compression, and selective headcount reduction rather than through any single aggregate shock. HR teams in high-exposure industries should monitor entry-level retention carefully — the floor-raiser dynamic implies that the outside options of lower-wage workers in AI-exposed roles are improving faster than their wages.

\paragraph{Limitations and future work} Four limitations warrant acknowledgment. First, the trend-break correction absorbs linear pre-trends but cannot address non-linear confounders. Second, the AEI treatment variable for the firm-level analysis was constructed from November 2025 data applied to a 2019--2024 panel — if occupational AI penetration rankings shifted over this period, our estimates are attenuated lower bounds. Third, Compustat covers only large publicly traded firms and may not generalize to the SME sector. Fourth, our post-period extends only to 2024; as AI capabilities accelerate, the effects documented here may grow substantially. Future work should exploit time-varying AEI data as it accumulates, disaggregate by industry (finance versus healthcare versus manufacturing), and examine hours worked and task composition beyond headcount.

% ── References ───────────────────────────────────────────────────────────────
\bibliographystyle{apalike}
\begin{thebibliography}{99}

\bibitem[Acemoglu \& Restrepo(2018)]{acemoglu2018}
Acemoglu, D., \& Restrepo, P. (2018).
The race between man and machine.
\textit{American Economic Review}, 108(6), 1488--1542.

\bibitem[Acemoglu et al.(2022)]{acemoglu2022}
Acemoglu, D., Autor, D., Hazell, J., \& Restrepo, P. (2022).
Artificial intelligence and jobs: Evidence from online vacancies.
\textit{Journal of Labor Economics}, 40(S1), S293--S340.

\bibitem[Autor et al.(2003)]{alm2003}
Autor, D.~H., Levy, F., \& Murnane, R.~J. (2003).
The skill content of recent technological change.
\textit{Quarterly Journal of Economics}, 118(4), 1279--1333.

\bibitem[Brynjolfsson et al.(2025)]{brynjolfsson2025}
Brynjolfsson, E., Li, D., \& Raymond, L. (2025).
Generative AI at work.
\textit{Quarterly Journal of Economics}, 140(2), 889--942.

\bibitem[Callaway \& Sant'Anna(2021)]{callaway2021}
Callaway, B., \& Sant'Anna, P.~H. (2021).
Difference-in-differences with multiple time periods.
\textit{Journal of Econometrics}, 225(2), 200--230.

\bibitem[Eloundou et al.(2024)]{eloundou2024}
Eloundou, T., Manning, S., Mishkin, P., \& Rock, D. (2024).
GPTs are GPTs: An early look at the labor market impact potential of large language models.
\textit{Science}, 384(6702), 1306--1308.

\bibitem[Felten et al.(2021)]{felten2021}
Felten, E., Raj, M., \& Seamans, R. (2021).
Occupational, industry, and geographic exposure to artificial intelligence.
\textit{Strategic Management Journal}, 42(12), 2195--2217.

\bibitem[Frey \& Osborne(2013)]{fo2013}
Frey, C.~B., \& Osborne, M.~A. (2013).
The future of employment: How susceptible are jobs to computerisation?
\textit{Oxford Martin School Working Paper}.

\bibitem[Hui et al.(2024)]{hui2024}
Hui, X., Reshef, O., \& Zhou, L. (2024).
The short-term effects of generative AI on employment: Evidence from an online labor market.
\textit{Organization Science}, 35(6), 1977--1989.

\bibitem[Massenkoff \& McCrory(2026)]{massenkoff2026}
Massenkoff, M., \& McCrory, E. (2026).
The Anthropic Economic Index.
\textit{Anthropic Working Paper}.

\bibitem[Noy \& Zhang(2023)]{noy2023}
Noy, S., \& Zhang, W. (2023).
Experimental evidence on the productivity effects of generative artificial intelligence.
\textit{Science}, 381(6654), 187--192.

\bibitem[Roth et al.(2023)]{roth2023}
Roth, J., Sant'Anna, P.~H., Bilinski, A., \& Poe, J. (2023).
What's trending in difference-in-differences?
\textit{Journal of Econometrics}, 235(2), 2218--2244.

\bibitem[Sun \& Abraham(2021)]{sun2021}
Sun, L., \& Abraham, S. (2021).
Estimating dynamic treatment effects in event studies with heterogeneous treatment effects.
\textit{Journal of Econometrics}, 225(2), 175--199.

\bibitem[Webb(2020)]{webb2020}
Webb, M. (2020).
The impact of artificial intelligence on the labor market.
\textit{Stanford University Working Paper}.

\bibitem[Wolfers(2006)]{wolfers2006}
Wolfers, J. (2006).
Did unilateral divorce laws raise divorce rates? A reconciliation and new results.
\textit{American Economic Review}, 96(5), 1802--1820.

\end{thebibliography}

\end{document}
